\section{An exact multiplicity-free rectangular family}

\subsection*{F-019. The full square of a two-row rectangle}
This is an extracted form of a known positive formula, not a
claim of a new Kronecker rule. Let $h\geq1$ and
$\nu\vdash2h$. If $\ell(\nu)\leq4$, pad $\nu$ by zeros
to four parts. Then
\[
 g(\nu,(h,h),(h,h))=
 \begin{cases}
 1,&\nu_1,\nu_2,\nu_3,\nu_4\text{ have the same parity},\\
 0,&\text{otherwise}.
 \end{cases}                                          \tag{F-019.1}
\]
The coefficient is zero if $\ell(\nu)>4$. Thus the positive
cases are exactly all-even row lengths (with zero parts
allowed), and four positive odd row lengths. This includes
every output shape, rather than only two-row outputs.

\paragraph{Known source and short extraction.}
The exact supplied source is Blasiak--Mulmuley--Sohoni,
GCT IV, Proposition 18.5, label
\texttt{p Kronecker consequence l1 = 0}, in
\texttt{Apr13KroneckerGCT4.tex}, lines 7806--7813. The
equivalent positive sum immediately before it, label
\texttt{e BWZ main result}, is at lines 7778--7784; the
paper attributes that sum to its reference BWZ. The column
notation is defined at lines 1535--1536 and used at lines
7453--7466. Numbering is with the
shared theorem counter in the supplied version; a current
attempt to reopen the author-hosted PDF was unavailable,
so no new PDF page verification is claimed here.

Write $c_i=\nu_i-\nu_{i+1}$, $1\leq i\leq3$.
In that source $(h,h)$ is $[h,0]$, and the coefficient of
$x^k$ in its generating function after replacing $x$ by
$x^{1/2}$ corresponds to the second two-row factor having
row difference $2k$. For a second rectangle, $k=0$.
The positive sum at that coefficient has just the term
$j=0$ in its first sum; the second sum is empty. That term
is one exactly when $c_1,c_2,c_3$ are all even, and zero
otherwise. The three differences are even precisely when
all four padded row lengths have the same parity. This
proves (F-019.1) directly from the stated known formula.

For the length assertion, realize the coefficient by
Schur--Weyl as the multiplicity of
$S_{(h,h)}\mathbb C^2\otimes S_{(h,h)}\mathbb C^2$
in $S_\nu(\mathbb C^2\otimes\mathbb C^2)$.
The latter is zero when $\ell(\nu)>4$.
This uses the characteristic-zero Schur--Weyl equivalence,
Macdonald, I, Appendix A, (5.3)--(5.6), pp.~155--156.

\paragraph{Complete stretching in this family.}
For every positive integer $t$,
\[
 g(t\nu,(th,th),(th,th))=
 \begin{cases}
 1,&t\text{ even, or the four parts of }\nu
                              \text{ already have one parity},\\
 0,&t\text{ odd and those parities are mixed},
 \end{cases}                                          \tag{F-019.2}
\]
provided $\ell(\nu)\leq4$; for greater length every stretch
is zero. This follows by multiplying all three adjacent
row differences by $t$ in (F-019.1). Thus every mixed-parity
shape of length at most four is an ordinary saturation hole,
and stretch two has coefficient exactly one. The conclusion
is consistent with, and a special exact-value case of, O-001.

\paragraph{Grothendieck dictionary and GCT scope.}
Let $G$ be the constant group scheme $S_{2h}$ over
$\mathbb Z$, as in SGA 3, I, 4.1, p.~23. Over
$\mathbb Q$, the finite-group invariant space
\[
   \bigl([\nu]\otimes[(h,h)]\otimes[(h,h)]\bigr)^G
\]
has rank exactly zero or one as specified above.
The Specht modules may be constructed by the finite-set
incidence kernels in A-018. Their tensor and Schur functors
are operations in the exact representation $K_0$, with
SGA 3, I, Proposition 4.7.2, pp.~34--35, and SGA 6, VI,
Theorem 3.3, pp.~392--393. No effective permutation lift
of an irreducible is asserted. This is an explicit positive
multiplicity rule in a rectangular subfamily; it does not
compare the evaluation kernels defining the two GCT orbit
closures and therefore supplies no additional actual
multiplicity obstruction on its own.
